\documentclass[11pt,a4paper]{ctexart}
\usepackage[margin=2.2cm]{geometry}
\usepackage{amsmath,amssymb,amsthm,mathtools,booktabs,hyperref,microtype}
\hypersetup{colorlinks=true,linkcolor=black,urlcolor=blue,citecolor=black}
\newtheorem{theorem}{定理}
\newtheorem{lemma}{引理}
\newtheorem{corollary}{推论}
\renewcommand{\proofname}{证明}
\newcommand{\Fthree}{F_3}
\newcommand{\vthree}{v_3}
\title{\bfseries 任意高阶的低次数刚性\\与总乘积深度分岔}
\author{研究札记 · F3-HIER}
\date{2026 年 9 月 29 日}
\begin{document}
\maketitle
\begin{abstract}
设 \(F_3(m)=v_3(m+1)-v_3(m-1)\)。本文构造一族具有 \(q=3^s\) 个素数端点的固定伸缩形状，使所有总次数小于 \(q\) 的单项式 \(F_3\) 数据完全固定，包括重复因子；但全部 \(q\) 个端点的总乘积深度仍可取任意预先指定的高层。核心机制是：低次数赋值始终停在模 \(3^{k+s}\) 的精度以内，而次数 \(q\) 首次触及边界。最高阶行为进一步由偏移系数和 \(C=\sum c_i\bmod3\) 控制：\(3\mid C\) 时总深度锁死，\(3\nmid C\) 时出现唯一三进根分支。对相应全素数族，得到显式渐近计数与几何深度律；普通等差族与末点移动族甚至拥有完全相同的素数计数主系数，却有不同的最高阶深度分布。
\end{abstract}

\section{定义与主定理}
对 \(m>1\)，
\[
\Fthree(m)=\vthree(m+1)-\vthree(m-1).
\]
固定
\[
k,s\ge1,\quad q=3^s,\quad b=k+s,\quad M=3^b,
\]
以及
\[
W=\prod_{\substack{\ell\le q\\ \ell\ {\rm prime},\ \ell\ne3}}\ell.
\]
取
\[
(c_0,\ldots,c_{q-1})=(0,1,\ldots,q-2,q),
\quad p_i=n+MWc_i d.
\]
令 \(\mathcal Q_{k,s}(X)\) 为满足 \(X<n,d\le2X\)、\(F_3(n)=k\)，且
\[
d,p_0,\ldots,p_{q-1}
\]
全部为素数的参数集合。

定义
\[
\begin{aligned}
\kappa_{k,s}:={}&
\frac4{3^{k+2}}\left(\frac32\right)^{q+1}
\prod_{\substack{\ell\le q\\\ell\ne3}}
\left(\frac{\ell}{\ell-1}\right)^{q-1}\\
&\times
\prod_{\ell>q}
\frac{\ell^{q-1}(\ell-q)}{(\ell-1)^q},
\end{aligned}
\]
其中乘积均遍历素数。

\begin{theorem}[F3-HIER-1]
若 \(\mathbf e=(e_i)\) 为非负整数指数向量，且
\[
1\le u:=\sum_i e_i<q,
\]
则
\[
\boxed{
\Fthree\!\left(\prod_i p_i^{e_i}\right)
=(-1)^{u+1}\bigl(k+\vthree(u)\bigr).
}
\]
此外
\[
|\mathcal Q_{k,s}(X)|
=
(\kappa_{k,s}+o(1))
\frac{X^2}{(\log X)^{q+1}}.
\]
若
\[
D=\Fthree\!\left(\prod_{i=0}^{q-1}p_i\right),
\]
则对每个固定 \(R\ge k+s\)
\[
\#\{D=R\}
=
(\kappa_{k,s}w_{k+s}(R)+o(1))
\frac{X^2}{(\log X)^{q+1}},
\]
其中
\[
w_b(R)=
\begin{cases}
1/2,&R=b,\\
3^{-(R-b)},&R>b.
\end{cases}
\]
\end{theorem}

\section{幂赋值与低次数刚性}
\begin{lemma}
若 \(F_3(n)=k\ge1\)，则对任意 \(u\ge1\)
\[
\vthree\bigl(n^u-(-1)^u\bigr)=k+\vthree(u),
\]
并且
\[
F_3(n^u)=(-1)^{u+1}(k+\vthree(u)).
\]
\end{lemma}
\begin{proof}
令 \(U=-n\in1+3\mathbb Z_3\)，则 \(v_3(U-1)=k\)。写 \(u=3^ta\)，\(3\nmid a\)。对 \(V\equiv1\pmod3\)，
\[
V^3-1=(V-1)(V^2+V+1),
\]
而 \(v_3(V^2+V+1)=1\)，故每立方一次赋值恰增加 \(1\)。于是
\[
v_3(U^{3^t}-1)=k+t.
\]
再因 \(3\nmid a\)，几何和 \((X^a-1)/(X-1)\) 在 \(X\equiv1\pmod3\) 时为单位，得到结论。
\end{proof}

若所有端点满足
\[
p_i\equiv n\pmod{3^{k+s}},
\]
则对单项式 \(P=\prod p_i^{e_i}\)
\[
P\equiv n^u\pmod{3^{k+s}}.
\]
当 \(u<3^s\) 时
\[
\vthree(u)\le s-1,
\]
所以
\[
k+\vthree(u)<k+s.
\]
因此模 \(3^{k+s}\) 已足以固定该精确赋值，主定理的低次数公式成立。

次数 \(q=3^s\) 时
\[
k+\vthree(q)=k+s,
\]
恰好首次碰到精度边界。

\section{系数和控制最高阶}
更一般地取
\[
p_i=n+MWc_i d,\qquad C=\sum_i c_i,
\]
其中 \(3\nmid W\)。由
\[
v_3(n^q+1)=k+s
\]
定义
\[
A(n)=\frac{n^q+1}{M}\in\mathbb Z_3^\times.
\]
展开总乘积可写成
\[
\prod_i p_i+1=M G_n(d),
\]
其中
\[
\boxed{
G_n(d)=A(n)+WCn^{q-1}d+M H_n(d).
}
\]
由于 \(q\) 为奇数，
\[
\Fthree\!\left(\prod_i p_i\right)=k+s+v_3(G_n(d)).
\]

若 \(3\mid C\)，则
\[
G_n(d)\equiv A(n)\not\equiv0\pmod3,
\]
所以总深度恒为 \(k+s\)。

若 \(3\nmid C\)，则
\[
G_n(d)\equiv A(n)+WCd\pmod3
\]
且
\[
G_n(d)-G_n(e)
=(d-e)\bigl(WCn^{q-1}+MK_n(d,e)\bigr).
\]
括号恒为单位，因此
\[
v_3(G_n(d)-G_n(e))=v_3(d-e).
\]
于是 \(G_n\) 是 \(\mathbb Z_3\) 上等距双射，存在唯一单位根 \(d_*(n)\)，且
\[
\boxed{
\Fthree\!\left(\prod_i p_i\right)
=
k+s+v_3(d-d_*(n)).
}
\]

普通系数 \(0,1,\ldots,q-1\) 的和被 \(3\) 整除；把最后一个点从 \(q-1\) 移到 \(q\)，系数和增加 \(1\)，因此从“锁死”切换为“开放根分支”，而低次数数据完全不变。

\section{全素数计数}
主形状取
\[
W=\prod_{\substack{\ell\le q\\\ell\ne3}}\ell.
\]
需要同时为素数的形式是
\[
d,\quad n+MWc_i d.
\]
方向 \((0,1)\) 与 \((1,MWc_i)\) 两两不成比例，系统有限复杂。

模 \(3\) 局部因子为
\[
\beta_3=\left(\frac32\right)^{q+1}.
\]
对 \(\ell\le q,\ell\ne3\)，因 \(\ell\mid W\)，所有偏移归零，仅需 \(n,d\) 为单位，得到
\[
\beta_\ell=\left(\frac{\ell}{\ell-1}\right)^{q-1}.
\]
对 \(\ell>q\)，\(c_i\) 模 \(\ell\) 两两不同，固定 \(d\ne0\) 后 \(n/d\) 需避开 \(q\) 个类，所以
\[
\beta_\ell=
\frac{\ell^{q-1}(\ell-q)}{(\ell-1)^q}
=
1+O_q(\ell^{-2}).
\]
因此奇异乘积严格为正。

对一个固定允许的模 \(3^A\) 参数类，有限复杂度素数线性形式定理给出
\[
\left(\frac{\mathfrak S_q}{3^{2A}}+o(1)\right)
\frac{X^2}{(\log X)^{q+1}}.
\]
满足 \(F_3(n)=k\) 的 \(n\) 类有
\[
2\cdot3^{A-k-1},
\]
单位 \(d\) 类有
\[
2\cdot3^{A-1}.
\]
求和得到
\[
\kappa_{k,s}=\frac4{3^{k+2}}\mathfrak S_q.
\]

唯一根分支在单位 \(d\) 中给出
\[
\Pr(v_3(G_n(d))=0)=1/2,
\qquad
\Pr(v_3(G_n(d))=t)=3^{-t}\quad(t\ge1).
\]
固定深度事件由有限模类决定，而每个允许类具有相同素数主系数，因此该比例原样转移到全素数集合。

\section{同主系数对照族}
比较
\[
(0,1,\ldots,q-1)
\]
与
\[
(0,1,\ldots,q-2,q).
\]
对于每个 \(\ell\le q,\ell\ne3\)，两族偏移都被 \(\ell\) 整除；对于 \(\ell>q\)，两族都给出 \(q\) 个互异排除类。因此每个素数处的局部因子完全相同，全素数主系数也相同。

然而普通族总深度恒为 \(k+s\)，移动族具有无限几何尾。这证明“相同低次数数据 + 相同素数计数主项”仍不足以决定最高阶深度分布。

\section{九端点与有限次数截断}
当 \(k=1,s=2\) 时
\[
q=9,\quad W=70,\quad M=27,
\]
所以
\[
p_i=n+1890c_i d,\qquad
(c_i)=(0,1,2,3,4,5,6,7,9).
\]
总次数 \(1,\ldots,8\) 的固定值为
\[
1,-1,2,-1,1,-2,1,-1.
\]
非空真子集有 \(510\) 个，含重复因子的低次数单项式共有
\[
\binom{17}{8}-1=24309.
\]

更一般地，给定任意次数上限 \(L\)，取 \(3^s>L\)，就得到全部次数 \(\le L\) 的 \(F_3\) 数据固定、而总乘积深度无界的全素数族。

\section{边界}
本文的全素数结论使用已有有限复杂度线性形式理论。\(q=3\) 时 Green--Tao 2010 的复杂度至多 \(2\) 无条件结果已经够用；一般 \(q\) 使用后续完成的一般有限复杂度版本。

固定 \(d\) 后系统退化为一变量仿射相关问题，本文不推出固定间距素数簇结论。

本稿状态为 \texttt{PAPER-AUDITED}，尚未 Lean 形式化。

\begin{thebibliography}{9}
\bibitem{GT10} Ben Green and Terence Tao, \emph{Linear equations in primes}, Ann. of Math. 171 (2010), 1753--1850.
\bibitem{GT12} Ben Green and Terence Tao, \emph{The Möbius function is strongly orthogonal to nilsequences}, Ann. of Math. 175 (2012), 541--566.
\bibitem{GTZ12} Ben Green, Terence Tao and Tamar Ziegler, \emph{An inverse theorem for the Gowers \(U^{s+1}[N]\)-norm}, Ann. of Math. 176 (2012), 1231--1372.
\end{thebibliography}
\end{document}
